\lampiransection{Template Triangulasi Data}
\label{app:template-triangulasi-data}

Lampiran ini digunakan untuk memeriksa konsistensi temuan penelitian melalui triangulasi sumber, triangulasi teknik, dan triangulasi waktu. Temuan dari observasi GrabFood sudah dapat dipakai sebagai data awal, tetapi keputusan strategis tetap perlu dikonfirmasi melalui wawancara, survei, dan dokumentasi.

\begin{table}[!htbp]
    \centering
    \caption{Triangulasi Awal Temuan Observasi GrabFood}
    \label{tab:triangulasi-awal-temuan-observasi-grabfood}
    \scriptsize
    \setlength{\tabcolsep}{3pt}
    \renewcommand{\arraystretch}{0.92}
    \begin{tabularx}{\textwidth}{|Z{2cm}|Y|Y|Y|Y|}
        \toprule
        \textbf{Kode} & \textbf{Temuan yang Diuji} & \textbf{Data Pendukung Saat Ini} & \textbf{Data Pembanding yang Dibutuhkan} & \textbf{Keputusan Sementara} \\ \hline
        \midrule
        TO-GF-01 & Reputasi digital Nasi Gerilya kuat dari rating 4,7 dan 4.584 penilaian. & Observasi profil GrabFood 13 Juni 2026. & Wawancara pelanggan, pemilik, dan observasi ulang rating. & Valid sebagai data observasi; perlu data tambahan untuk makna reputasi. \\ \hline
        TO-GF-03 & Promo/voucher aktif menjadi bagian penting dari daya tarik platform. & 19 promo/voucher terlihat pada observasi. & Wawancara pemilik tentang biaya promo dan survei pelanggan tentang pengaruh promo. & Perlu data tambahan. \\ \hline
        TO-GF-04 & Menu terlaris dapat menjadi fokus produk unggulan. & Label terlaris pada tiga menu utama. & Data penjualan/menu terlaris dari pemilik dan preferensi pelanggan. & Valid dengan catatan. \\ \hline
        TO-GF-05 & Rentang harga Nasi Gerilya lebar dan perlu dibaca bersama ketersediaan item. & Observasi rentang harga Nasi Gerilya dan kompetitor. & Wawancara pemilik tentang margin; survei pelanggan tentang persepsi nilai. & Valid dengan catatan. \\ \hline
        TO-GF-06 & Menu tidak tersedia menunjukkan potensi isu stok atau pembaruan aplikasi. & 6 menu/produk terlihat tidak tersedia. & Wawancara admin/karyawan dan observasi ulang pada jam berbeda. & Perlu data tambahan. \\ \hline
        TO-GF-07 & Rasa, porsi, dan kemasan menjadi kekuatan pengalaman pelanggan. & Ulasan positif dan ringkasan ulasan agregat. & Wawancara pelanggan dan survei pendukung. & Valid dengan catatan bias ulasan. \\ \hline
        TO-GF-08 & Akurasi pesanan, kelengkapan item, add-on, dan konsistensi kualitas menjadi area perbaikan. & Ulasan negatif/netral yang terlihat di GrabFood. & Wawancara admin/karyawan, pelanggan, dan catatan komplain. & Perlu data tambahan. \\ \hline
        \begin{tabular}[t]{@{}l@{}}TO-KP-01--\\TO-KP-07\end{tabular} & Kompetitor memberi ancaman pada rating, jumlah penilaian, jarak, harga, penanda reputasi, dan persepsi nilai. & Observasi Paripurna, Dendeng Batokok, Istana Krakatau, Pondok Gurih, dan Garuda tanggal 19 dan 20 Juni 2026. & Wawancara pemilik/pelanggan dan observasi kompetitor lanjutan pada waktu berbeda. & Valid sebagai pembanding awal, tetapi perlu konfirmasi relevansi kompetitor dari sisi pelanggan dan pemilik. \\ \hline
        \bottomrule
    \end{tabularx}
    \tablesource{Diolah peneliti berdasarkan observasi GrabFood dan rencana triangulasi data penelitian utama (2026).}
\end{table}

\begin{table}[!htbp]
    \centering
    \caption{Mitigasi Bias pada Data Ulasan GrabFood}
    \label{tab:mitigasi-bias-ulasan-grabfood}
    \footnotesize
    \setlength{\tabcolsep}{3pt}
    \renewcommand{\arraystretch}{0.95}
    \begin{tabularx}{\textwidth}{|Z{3cm}|Y|Y|}
        \toprule
        \textbf{Potensi Bias} & \textbf{Dampak pada Analisis} & \textbf{Cara Menanggulangi} \\ \hline
        \midrule
        \textit{Survivorship bias} ulasan & Ulasan hanya berasal dari pelanggan yang menyelesaikan pesanan dan memilih memberi review; pelanggan diam, batal pesan, atau berhenti membeli tidak terlihat. & Tambahkan pertanyaan wawancara dan survei tentang pengalaman yang tidak ditulis di GrabFood, alasan tidak mengulas, dan alasan berhenti membeli. \\ \hline
        Urutan dan penyaringan platform & Platform dapat menonjolkan ulasan terbaru, ulasan tertulis, atau topik tertentu sehingga contoh yang terlihat belum tentu mewakili seluruh pelanggan. & Baca ulasan positif, negatif, dan netral secara eksplisit; catat waktu pengambilan data dan lakukan observasi ulang. \\ \hline
        Bias pelanggan ekstrem & Pelanggan yang sangat puas atau sangat kecewa lebih mungkin meninggalkan ulasan. & Bandingkan ulasan dengan wawancara pelanggan biasa, pelanggan ulang, dan pelanggan yang jarang memberi review. \\ \hline
        Bias waktu observasi & Status toko, promo, menu tidak tersedia, dan rating dapat berubah sesuai waktu. & Lakukan observasi lanjutan pada hari kerja, akhir pekan, jam makan siang, dan malam hari. \\ \hline
        \bottomrule
    \end{tabularx}
    \tablesource{Diolah peneliti sebagai panduan triangulasi data ulasan GrabFood (2026).}
\end{table}

\begin{table}[!htbp]
    \centering
    \caption{Kriteria Keputusan Triangulasi}
    \label{tab:kriteria-keputusan-triangulasi}
    \footnotesize
    \setlength{\tabcolsep}{4pt}
    \renewcommand{\arraystretch}{0.95}
    \begin{tabularx}{\textwidth}{|Z{3cm}|Y|}
        \toprule
        \textbf{Keputusan} & \textbf{Kriteria} \\ \hline
        \midrule
        Valid & Temuan didukung oleh minimal dua sumber data atau dua teknik pengumpulan data dan tidak bertentangan dengan data waktu lain. \\ \hline
        Valid dengan catatan & Temuan didukung data, tetapi masih memiliki keterbatasan seperti jumlah informan terbatas, bias ulasan, atau data dokumentasi tidak lengkap. \\ \hline
        Perlu data tambahan & Temuan hanya berasal dari satu sumber atau masih perlu dikonfirmasi melalui wawancara, survei, dokumentasi, atau observasi waktu berbeda. \\ \hline
        Tidak digunakan & Temuan tidak memiliki dukungan data yang cukup atau tidak relevan dengan fokus penelitian. \\ \hline
        \bottomrule
    \end{tabularx}
    \tablesource{Diolah peneliti sebagai kriteria validasi temuan (2026).}
\end{table}
